% =============================================================
%  Chapter — DC Power Transmission (Module 03)
%  Figures in ./fig  (EMF & SVG converted to cropped PDF, PNG/JPG copied).
%  Note: figures are referenced WITH extension because some basenames
%        exist as both .pdf and .png.
% =============================================================
\newcommand{\dcfig}[1]{chapters/hvdc-topologies/fig/#1}

\sectionslide{DC Power Transmission}{Converters, topologies and load flow}

% ---- Converter ----------------------------------------------
\begin{frame}{Converter}
  \centering
  \topoimg[width=0.7\linewidth,height=0.3\textheight,keepaspectratio]{\dcfig{image4.pdf}}\\[0.6em]
  \begin{columns}[c,onlytextwidth]
    \begin{column}{0.48\textwidth}\centering
      \textbf{Inverter}\\[0.3em]
      \topoimg[width=\linewidth,height=0.3\textheight,keepaspectratio]{\dcfig{image5.pdf}}
    \end{column}
    \begin{column}{0.48\textwidth}\centering
      \textbf{Rectifier}\\[0.3em]
      \topoimg[width=\linewidth,height=0.3\textheight,keepaspectratio]{\dcfig{image6.pdf}}
    \end{column}
  \end{columns}
\end{frame}

\begin{frame}{Converter - Load flow definition}
  \begin{columns}[c,onlytextwidth]
    \begin{column}{0.56\textwidth}
      \begin{itemize}
        \setlength{\itemsep}{0.4em}
        \item Positive load is defined for the inverter case (generator
              convention).
        \item In the USA a common term is ``Inverter-Based Resources'' (IBR).
        \item IBR are not only HVDC systems but all kinds of generators:
          \begin{itemize}
            \item Wind Turbine Generators (WTG)
            \item PV, \ldots
          \end{itemize}
        \item A VSC-HVDC system can be an inverter or a rectifier depending
              on the power-flow direction; some applications require changing
              direction.
      \end{itemize}
    \end{column}
    \begin{column}{0.40\textwidth}\centering
      \topoimg[width=\linewidth,height=0.6\textheight,keepaspectratio]{\dcfig{image7.pdf}}
    \end{column}
  \end{columns}
\end{frame}

\begin{frame}{Converter terminology}
  \centering
  \topoimg[width=\linewidth,height=0.72\textheight,keepaspectratio]{\dcfig{image8.png}}\\[0.3em]
  {\color{atGray}\footnotesize IEC 62747:2014}
\end{frame}

\begin{frame}{DC load flow definitions}
  \begin{columns}[c,onlytextwidth]
    \begin{column}{0.4\textwidth}\centering
      \begin{itemize}
        \item DC equipment (incl. DC line) voltage is defined by the voltages $U_{dpe\_rec}$ and $U_{dpe\_inv}$
        \item Converter size is defined by $U_{dpp\_rec}$
        \item HVDC systems have bidirectional transmission capability, which means that every individual converter can be either rectifier or inverter
      \end{itemize}
    \end{column}
    \begin{column}{0.6\textwidth}\centering
      \topoimg[width=\linewidth,height=0.5\textheight,keepaspectratio]{\dcfig{image9.pdf}}
    \end{column}
  \end{columns}
\end{frame}

% ---- Monopole -----------------------------------------------
\begin{frame}{Monopole}
  \begin{columns}[c,onlytextwidth]
    \begin{column}{0.6\textwidth}\centering
      \topoimg[width=\linewidth,height=0.72\textheight,keepaspectratio]{\dcfig{image12.pdf}}
    \end{column}
    \begin{column}{0.40\textwidth}
      \textbf{DC load flow}
      \begin{align*}
        U_{dpp\_rec} &= U_{dpe\_rec}\\
        I_{DC} &= \tfrac{U_{dpp\_rec}-U_{dpp\_inv}}{2 \cdot R_{DC}}\\
        U_{dne\_rec} &= I_{d}\cdot R_{DC}
      \end{align*}
      \begin{itemize}
        \item Max.\ voltage created by converter: $U_{dpp\_rec}$
        \item Max. equipment voltage: 
        \item \textcolor{red}{DC cables have different voltage ratings}
      \end{itemize}
    \end{column}
  \end{columns}
\end{frame}

\begin{frame}{Monopole - extending power transmission}
  \begin{columns}[c,onlytextwidth]
    \begin{column}{0.4\textwidth}
      \begin{itemize}
        \setlength{\itemsep}{0.3em}
        \item \textbf{Higher DC voltage}
            \begin{itemize}
              \item Limited equipment availability
              \item No redundancy
            \end{itemize}
        \item \textbf{Higher DC current}
            \begin{itemize}
              \item Higher DC line costs
              \item No redundancy
            \end{itemize}
        \item \textbf{Additional monopole}
            \begin{itemize}
              \item Additional converter equipment
              \item 4 DC conductors
              \item Increased reliability
            \end{itemize}
      \end{itemize}
    \end{column}
    \begin{column}{0.6\textwidth}\centering
      \topoimg[width=\linewidth,height=0.34\textheight,keepaspectratio]{\dcfig{image14.pdf}}\\[0.4em]
      \topoimg[width=\linewidth,height=0.34\textheight,keepaspectratio]{\dcfig{image13.pdf}}
    \end{column}
  \end{columns}
\end{frame}

\begin{frame}{Monopole $\rightarrow$ Bipole (BIP)}
  \centering
  \begin{columns}[c,onlytextwidth]
    \begin{column}{0.42\textwidth}\centering
      \topoimg[width=\linewidth,height=0.5\textheight,keepaspectratio]{\dcfig{image13.pdf}}\\[0.4em]
      \topoimg[width=\linewidth,height=0.5\textheight,keepaspectratio]{\dcfig{image14.pdf}}
    \end{column}
    \begin{column}{0.06\textwidth}\centering $\Rightarrow$ \end{column}
    \begin{column}{0.42\textwidth}\centering
      \topoimg[width=\linewidth,height=0.8\textheight,keepaspectratio]{\dcfig{image15.pdf}}
    \end{column}
  \end{columns}
\end{frame}

\begin{frame}{Bipole (BIP)}
  \begin{columns}[c,onlytextwidth]
    \begin{column}{0.54\textwidth}
      \begin{itemize}
        \setlength{\itemsep}{0.4em}
        \item 3 conductors of 2 different types.
        \item Two poles that can transmit power independently.
        \item On a fault in one pole, the other pole still transmits power ---
              redundancy.
        \item In symmetrical operation (both poles equal), no current flows
              through the DMR.
      \end{itemize}
    \end{column}
    \begin{column}{0.42\textwidth}\centering
      \topoimg[width=\linewidth,height=0.66\textheight,keepaspectratio]{\dcfig{image16.pdf}}
    \end{column}
  \end{columns}
\end{frame}

\begin{frame}{BIP operation modes}
  \begin{columns}[T,onlytextwidth]
    \begin{column}{0.32\textwidth}\centering
      \topoimg[width=\linewidth,height=0.42\textheight,keepaspectratio]{\dcfig{image17.pdf}}\\[0.3em]
      {\footnotesize\textbf{Bipolar --- balanced}\\ 
      \begin{itemize}
        \item Equal power for poles
        \item No DMR current
        \item Loss-optimized
      \end{itemize}
      }
    \end{column}
    \begin{column}{0.32\textwidth}\centering
      \topoimg[width=\linewidth,height=0.42\textheight,keepaspectratio]{\dcfig{image18.pdf}}\\[0.3em]
      {\footnotesize\textbf{Monopolar}\\
      \begin{itemize}
        \item Emergency mode
        \item System-design definition
      \end{itemize}
      }
    \end{column}
    \begin{column}{0.32\textwidth}\centering
      \topoimg[width=\linewidth,height=0.42\textheight,keepaspectratio]{\dcfig{image19.pdf}}\\[0.3em]
      {\footnotesize\textbf{Bipolar --- unbalanced}\\
      \begin{itemize}
        \item Split-busbar operation
        \item DMR current flow
      \end{itemize}
      }
    \end{column}
  \end{columns}
\end{frame}

\begin{frame}{DC load flow}
  \begin{columns}[c,onlytextwidth]
    \begin{column}{0.5\textwidth}\centering
      \topoimg[width=\linewidth,height=0.4\textheight,keepaspectratio]{\dcfig{image20.pdf}}
      
    \end{column}
    \begin{column}{0.50\textwidth}
      {\centering
      \topoimg[width=\linewidth,height=0.4\textheight,keepaspectratio]{\dcfig{image21.pdf}}\\[0.4em]}

      For DC load-flow calculation the converter can be represented as an ideal voltage source.
    \end{column}
  \end{columns}
\end{frame}

\begin{frame}{Balanced operation}
  \centering
  \topoimg[width=\linewidth,height=0.74\textheight,keepaspectratio]{\dcfig{image22.png}}
\end{frame}

\begin{frame}{Balanced operation - power definition}
  \begin{columns}[c,onlytextwidth]
    \begin{column}{0.48\textwidth}\centering
      \topoimg[width=\linewidth,height=0.66\textheight,keepaspectratio]{\dcfig{image22.png}}
    \end{column}
    \begin{column}{0.48\textwidth}\centering
      \topoimg[width=\linewidth,height=0.66\textheight,keepaspectratio]{\dcfig{image24.png}}
    \end{column}
  \end{columns}
\end{frame}

\begin{frame}{Balanced vs. unbalanced operation}
  \begin{columns}[c,onlytextwidth]
    \begin{column}{0.48\textwidth}\centering
      \textbf{Balanced mode}\\[0.3em]
      \topoimg[width=\linewidth,height=0.6\textheight,keepaspectratio]{\dcfig{image22.png}}
    \end{column}
    \begin{column}{0.48\textwidth}\centering
      \textbf{Unbalanced mode}\\[0.3em]
      \topoimg[width=\linewidth,height=0.6\textheight,keepaspectratio]{\dcfig{image26.png}}
    \end{column}
  \end{columns}
\end{frame}

\begin{frame}{Monopolar operation}
  \begin{columns}[c,onlytextwidth]
    \begin{column}{0.6\textwidth}\centering
      \topoimg[width=\linewidth,height=0.7\textheight,keepaspectratio]{\dcfig{image28.png}}
    \end{column}
    \begin{column}{0.40\textwidth}
      \begin{itemize}
        \setlength{\itemsep}{0.4em}
        \item Usually for a short period, until system reconfiguration is
              finished.
        \item For longer operation, the DMR must be designed accordingly.
      \end{itemize}
    \end{column}
  \end{columns}
\end{frame}

\begin{frame}{System earthing}
  \begin{columns}[c,onlytextwidth]
    \begin{column}{0.48\textwidth}\centering
      \topoimg[width=\linewidth,height=0.62\textheight,keepaspectratio]{\dcfig{image28.png}}
    \end{column}
    \begin{column}{0.48\textwidth}\centering
      \topoimg[width=\linewidth,height=0.62\textheight,keepaspectratio]{\dcfig{image30.png}}
    \end{column}
  \end{columns}
  \vspace{0.3em}
  Earthing of the BIP system has no impact on the resulting currents and losses.
\end{frame}

% ---- Rigid Bipole -------------------------------------------
\begin{frame}{BIP $\rightarrow$ Rigid BIP}
  \begin{columns}[c,onlytextwidth]
    \begin{column}{0.48\textwidth}\centering
      \textbf{Bipole}\\[0.3em]
      \topoimg[width=\linewidth,height=0.62\textheight,keepaspectratio]{\dcfig{image16.pdf}}
    \end{column}
    \begin{column}{0.48\textwidth}\centering
      \textbf{Rigid Bipole}\\[0.3em]
      \topoimg[width=\linewidth,height=0.62\textheight,keepaspectratio]{\dcfig{image32.pdf}}
    \end{column}
  \end{columns}
\end{frame}

\begin{frame}{Rigid BIP}
  \begin{columns}[c,onlytextwidth]
    \begin{column}{0.54\textwidth}
      \begin{itemize}
        \setlength{\itemsep}{0.4em}
        \item 2 conductors of 1 type.
        \item Both poles must transmit the same power (same current flows
              through both poles).
        \item On a fault in one pole, the other pole might not be able to
              transmit the power.
      \end{itemize}
    \end{column}
    \begin{column}{0.42\textwidth}\centering
      \topoimg[width=\linewidth,height=0.62\textheight,keepaspectratio]{\dcfig{image32.pdf}}
    \end{column}
  \end{columns}
\end{frame}

\begin{frame}{Rigid BIP operation modes}
  \begin{columns}[T,onlytextwidth]
    \begin{column}{0.32\textwidth}\centering
      \topoimg[width=\linewidth,height=0.42\textheight,keepaspectratio]{\dcfig{image33.pdf}}\\[0.3em]
      {\footnotesize\textbf{Common-busbar}\\ 
      \begin{itemize}
        \item Standard operation mode.
      \end{itemize}
      }
    \end{column}
    \begin{column}{0.32\textwidth}\centering
      \topoimg[width=\linewidth,height=0.42\textheight,keepaspectratio]{\dcfig{image34.pdf}}\\[0.3em]
      {\footnotesize\textbf{Split-busbar}\\
      \begin{itemize}
        \item Poles must keep equal power in all modes (incl.\ dynamic events)
      \end{itemize}
      }
    \end{column}
    \begin{column}{0.32\textwidth}\centering
      \topoimg[width=\linewidth,height=0.42\textheight,keepaspectratio]{\dcfig{image35.pdf}}\\[0.3em]
      {\footnotesize\textbf{Monopolar}\\
      \begin{itemize}
        \item Reconfiguration takes $>1$\,s.
      \end{itemize}
      }
    \end{column}
  \end{columns}
\end{frame}

% ---- Symmetrical Monopole -----------------------------------
\begin{frame}{Symmetrical Monopole (SMP)}
  \begin{columns}[c,onlytextwidth]
    \begin{column}{0.54\textwidth}\centering
      \topoimg[width=\linewidth,height=0.72\textheight,keepaspectratio]{\dcfig{image37.pdf}}
    \end{column}
    \begin{column}{0.42\textwidth}
      \begin{itemize}
        \setlength{\itemsep}{0.25em}
        \item High-impedance earth reference on the AC side
        \item $U_{dpp\_Rec}= 2 \cdot \lvert U_{dpe\_Rec} \rvert$
        \item $I_d=\frac{U_{dpp\_rec} - U_{dpp\_inv}}{2 \cdot R_d}$
        \item Only for VSC-HVDC systems
        \item \textcolor{green}{DC cables have the same voltage rating} 
      \end{itemize}
    \end{column}
  \end{columns}
\end{frame}

% ---- MMC comparison table -----------------------------------
\begin{frame}{Modular Multilevel Converter (MMC)}
  Power transmission parameters:
  \begin{itemize}
    \item Power: 1000\,MW
    \item DC voltage: $\pm$ 320 kV
    \item Max. DC current: 1563 A
  \end{itemize}
  \vspace{0.6em}

  \centering
  \renewcommand{\arraystretch}{1.3}
  \begin{tabular}{@{}l c c c@{}}
     & \textbf{2 $\times$ Monopole} & \textbf{1 $\times$ BIP} & \textbf{1 $\times$ SMP} \\
    \hline
    Transformers      & \textcolor{red}{$--$} & \textcolor{red}{$--$} & \textcolor{green}{$++$} \\
    DC lines          & \textcolor{red}{4 lines / 2 types} & \textcolor{red}{3 lines / 2 types} & \textcolor{green}{2 lines / 1 type}\\
    Converters        & \textcolor{red}{4 converters} & \textcolor{red}{4 converters} & \textcolor{green}{2 converters}\\
    Converter rating ($U_{dpp}$) & \textcolor{green}{320\,kV} & \textcolor{green}{320\,kV} & \textcolor{red}{640\,kV} \\
    Losses            & \textcolor{red}{$--$} & \textcolor{red}{$-$} & \textcolor{green}{$+$} \\
    Reliability       & \textcolor{red}{$++$} & \textcolor{red}{$+$} & \textcolor{red}{$--$} \\
    \hline
  \end{tabular}
\end{frame}

% ---- DC cable -----------------------------------------------
\begin{frame}{DC cable}
  \begin{columns}[c,onlytextwidth]
    \begin{column}{0.56\textwidth}
      \begin{itemize}
        \setlength{\itemsep}{0.25em}
        \item \textbf{Insulator} --- XLPE (cross-linked polyethylene) or
              MI (mass-impregnated).
        \item \textbf{Conductor} --- copper or aluminium.
        \item Balance between DC voltage and DC current for the most
              economical design.
        \item Insulator is generally cheaper than conductor
              $\rightarrow$ higher voltage rating is preferred
      \end{itemize}
      \vspace{0.3em}
      Project voltages: 80/200/250/\textbf{320}/400/\textbf{525}\,kV\\
      (525\,kV is new standard for $\ge 2000$\,MW).
    \end{column}
    \begin{column}{0.40\textwidth}\centering
      \topoimg[width=\linewidth,height=0.4\textheight,keepaspectratio]{\dcfig{image39.png}}\\[0.4em]
      \topoimg[width=0.8\linewidth,height=0.28\textheight,keepaspectratio]{\dcfig{image38.jpg}}\\
      {\color{atGray}\scriptsize www.borealisgroup.com}
    \end{column}
  \end{columns}
\end{frame}

% ---- Topology comparison table ------------------------------
\begin{frame}{Topology comparison}
  \centering
  \renewcommand{\arraystretch}{1.3}
  \begin{tabular}{@{}p{4.6cm} c c c@{}}
     & \textbf{Sym.\ Monopole} & \textbf{Rigid Bipole} & \textbf{Bipole w/ MR} \\
    \hline
    DC conductors                         & 2 & 2 & 3 \\
    Power on converter fault              & 0\% & 50\% & 50\% \\
    Power on DC conductor fault           & 0\% & 0\% & 50\% \\
    Economical DC voltage                 & $\le 400$\,kV & $>500$\,kV & $>500$\,kV \\
    Costs                                 & \$ & \$\$ & \$\$\$ \\
    Max.\ power                           & \makecell{1400\,MW (320\,kV)\\1500\,MW (400\,kV)} & 2000\,MW & 3000\,MW \\
    \hline
  \end{tabular}
\end{frame}
